\subsection{A continuation under a new root}
The central question is whether an old call still belongs to a valid
provisional continuation after trusted feedback advances. Its birth version
records dispatch history. The dependency witness establishes whether it
remains usable now.

Let $x=(F,P_C)$. A complete recorded node provides an edge
\begin{equation}
 (F,P_C)\xrightarrow{n}(b_n,P_n^+)
 \quad\text{if}\quad a_n=F\ \land\ P_n^-=P_C,
 \label{eq:edge}
\end{equation}
under the input, identity, and execution contract of
Section~\ref{sec:contract}. A finite sequence of such edges is a
\emph{valid provisional continuation} of $C$ when each successor consumes
the preceding complete output state at its actual frontier. Write
$C\leadsto n$ for a complete path to $n$. A frozen in-flight successor can
belong to the same continuation, but contributes no usable result until
completion.

\begin{proposition}[Sufficient re-rooting witness]
\label{prop:rebase}
Suppose an old branch contains a complete boundary node at the new certified
frontier $F$, with output state $P_C$. If its subsequent recorded calls form
an unbroken input--output path to the actual branch tip, those nodes are a
valid provisional continuation of $C$. An in-flight successor may continue
when its unchanged frozen input is that tip and its target exceeds $F$.
\end{proposition}
\begin{proof}
The boundary node provides $(F,P_C)$, equal to the new root's declared
cross-call state. The first successor therefore satisfies
Eq.~\eqref{eq:edge}. Inductively, each next call consumes the preceding
complete output at its recorded frontier, with the remaining inputs bound
by the same contract. The path reaches the actual tip. The in-flight
successor satisfies the next input edge, without yet supplying an output.
Re-rooting changes none of these calls or their frozen inputs.
\end{proof}

This establishes dependency validity. The provisional result remains subject
to later selection and revision; re-rooting grants no Canonical write
authority. The argument also makes no bitwise-replay assumption for a
different floating-point execution. Its witness follows calls that actually
occurred.

Two short counterexamples explain the required observations. A branch born
under $v$ may reach exactly $(F,P_C)$ when $v+1$ is installed; discarding it
solely because its version is older loses a valid continuation. Conversely,
equal displayed text can hide different feedback tokens or alignment state.
A call consuming one state has no valid path from the other merely because
their text agrees. Equal state at a different audio frontier likewise does
not provide Eq.~\eqref{eq:edge}.

\subsection{Retaining computation: \texorpdfstring{$R$}{R}}
In the core Qwen implementation, the computation guard implements a conservative witness for
Proposition~\ref{prop:rebase}. On a Canonical installation, it first finds a
complete Private node at exactly $F$ whose full feedback tokens equal
$P_C$. It follows complete successors, checking each frozen input and
frontier, until it reaches the recorded branch tip. An intermediate match
cannot stand in for a path to that tip. If a successor is in flight, the
guard requires the unique tracked task to consume this tip and target
audio beyond $F$.

When the witness succeeds, the Owner retains the branch and updates its
eligibility relative to the new root. The original job identifiers, birth
versions, input tokens, and results remain unchanged. Otherwise it fences
the continuation lacking this witness and rebuilds from an allowed state. This
guard's boundary-node and unique-slot requirements restrict what this
implementation can retain; they are not universal necessities for every
valid reuse scheme.

\subsection{Retaining an actual view: \texorpdfstring{$S$}{S}}
Computation retention alone does not preserve service that has already
reached the Consumer. A snapshot Consumer can replace an ahead view with
the newly certified, shorter view even when the ahead computation is still
valid. RevTrace instead maintains a certified lane and a provisional lane.
Both hold whole hypotheses; there is no uncertified suffix splice. The
Consumer records the certified text and provisional visible text separately,
so a retained provisional view acquires no certified-prefix status.

On applying a new Canonical publication, $S$ checks the complete currently
visible provisional node. It can prove a path directly when that node's
frozen input is $(F,P_C)$, or follow its complete published ancestors to
that boundary. Missing, cyclic, or ambiguous ancestry is rejected. This
check uses the same feedback-state contract as $R$, while accepting a
different sufficient witness: a directly rooted published call does not
need an extra Private boundary node.

The Consumer first previews the whole next view. The protocol kernel then
checks the installed durable root, actual registered results, dependency
path, and bound text/frontier before authorizing the actual application.
Only that application yields the recorded view receipt. A caller-supplied
compatibility flag cannot replace the underlying nodes.

\begin{proposition}[Continuity at a compatible application]
\label{prop:continuity}
Immediately before an authorized Canonical application, let the visible
whole provisional view $h$ be ahead of its new frontier $F$. If $S$ proves
$C\leadsto h$ and no independent withdrawal mode is invoked, that
application advances the certified lane and preserves $h$ unchanged.
\end{proposition}
\begin{proof}
The compatible-and-ahead branch installs $C$ into the certified lane and
retains the existing provisional node, text, and frontier. The kernel
authorizes this preview against the actual root and complete dependency
records before the Consumer applies it. Thus this application does not
itself withdraw any part of $h$.
\end{proof}

An incompatible update must still withdraw the view. Later valid results
can also revise text or replace a hypothesis. Consequently the property
concerns one compatible root application, rather than global monotonic
display. Snapshot-only service and conservative recomputation can be safe;
the four-arm experiment asks what additional progress $R$ and $S$ preserve
together, and where that benefit is worth its cost.
